\chapter{塑料、橡胶和纤维}

\begin{kaichang}
请你去屋里拿三样东西：一个喝空了的矿泉水瓶、一根橡皮筋，再找一件标签上写着“涤纶”或“聚酯纤维”的衣服——找不到的话，拿一个超市塑料袋也行。

现在猜两个问题。第一：这三样东西，一个硬、一个软、一个能拉成细丝，手感天差地别，可化学家却说它们是“同一类东西”，凭什么？第二：同样是“塑料”，为什么矿泉水瓶摔不烂、埋不烂，橡皮筋晒上几年却会发黏、开裂、一扯就断？

把猜测写下来。这一章结束时，你要能用“一根极长的线”把这三样东西全部解释清楚——包括为什么你的衣服标签上，“涤纶”和“聚酯纤维”后面往往跟着同一个英文缩写：PET。
\end{kaichang}

\section{又轻又韧，还不肯烂}

先只看能看见、能摸到的事实。

一个 500 毫升的玻璃瓶大约重三四百克，同样大小的塑料瓶只有二十克上下，还摔不碎。它透明、能弯、能拧，装满水放上一年，瓶壁既不溶解也不发霉。把它埋进土里，几十年后挖出来，形状还在——这是它的优点，也是它的麻烦。

橡皮筋恰恰相反：它软，能拉到原长的三四倍还不破，一松手“啪”地弹回去。可是把橡皮筋忘在窗台上晒一个夏天，它会变黏、发脆，轻轻一扯就断。

涤纶衣服轻、薄、干得快、几乎不起皱，可穿着它靠近烧烤架，一个火星溅上去，面料不是烧焦，而是熔出一个边缘发硬的小洞。

这三组现象——“韧而不烂”“能拉能弹却会老化”“遇火熔而不燃”——看起来互不相干。但它们的背后，是同一种粒子层面的结构：\textbf{极长的分子链}。

\section{小分子串成的长链}

回忆一下第 36 章：碳原子最多能伸出四只“手”（四个共价键），所以能连成链、环和网。前面几章遇到的有机分子——乙醇、醋酸、阿司匹林——都是几十个原子的小分子。而现在，请你想象一个碳原子牵起另一个，再牵一个，一直重复几千次。

原料是一种小得不能再小的分子：乙烯，两个碳、四个氢，两个碳之间有一双键（第 36 章讲过双键）。在化工厂里，成千上万个乙烯分子一个接一个地接上去，双键打开，变成一条由几千个碳原子组成的长链。乙烯叫\textbf{单体}（单独的一个小分子），接出来的长链叫\textbf{聚合物}（“聚”就是聚合，“合物”是化合物）——这条链的商品名字你很熟：聚乙烯，塑料袋和牛奶桶的材料。

一个聚乙烯分子有几千个碳原子，可它和邻居分子之间\textbf{没有}化学键相连——链与链之间只有第 22 章讲过的分子间作用力，弱而多。这就是理解一切塑料性质的钥匙：想象一大碗煮熟的面条。每根面条内部非常结实（碳—碳共价键，第 19 章说过，断它要花很大力气），但面条和面条之间可以滑动、可以缠结。用力拉这碗面，面条先是被捋直，然后彼此滑开——材料就变形了；加热让分子抖动加剧，滑动更容易——材料就软了、熔了。塑料的“可塑”，说的就是链与链之间的滑动，而不是链本身断开。

“不肯烂”也说得通了。碳—碳主链上没有任何一处是微生物的酶认得的“接口”，水也泡不动它。降解不是方向不允许，而是慢得看不到头——方向归热力学管，快慢归动力学管，这是第 29 章的老规矩。

\section{两种串法：接龙与牵手}

把小分子串成长链，化学家主要用两种办法，区别只在于“接口”。

\textbf{第一种：加成聚合，玩的是接龙。}带双键的单体（乙烯、氯乙烯、苯乙烯……）把双键打开，像一列人手拉手排队，一个接一个接上去，链上只留下单键。所有原料原子全部进入长链，不产生任何副产物。让这场接龙开始，需要有人“推第一把”：老式工艺用高温高压加自由基引发剂，做出来的是支链很多的低密度聚乙烯（LDPE），软而透明，就是塑料袋的手感；后来齐格勒和纳塔发明了特殊催化剂（第 30 章：催化剂只换反应路径，不改方向和平衡），在低温低压下让乙烯规规矩矩排成直链，得到高密度聚乙烯（HDPE），硬而乳白，是牛奶桶和下水管的材料。同一种单体，换了“开工方式”，材料判若两物——这个差别马上还会用到。

\textbf{第二种：缩合聚合，玩的是牵手。}每个单体分子带两个“反应接口”（第 37 章的官能团），两个接口相遇，连上一根新键，同时脱掉一个小分子——通常是水。第 40 章讲过羧酸和醇相遇生成酯、同时出水（酯化），现在把这个反应搬到“两头都有接口”的分子上：对苯二甲酸（两头都是羧基）遇上乙二醇（两头都是羟基），牵一次手生成一个酯键、脱掉一分子水，剩下的两头还能继续牵，于是链越来越长——产物就是聚酯，代号 PET，做饮料瓶，拉成丝就叫涤纶。把羧酸换成己二酸、醇换成两头是氨基的己二胺，牵手生成的是酰胺键，链就叫尼龙。

记住一条账本规则：加成聚合不欠账，多少原子进、多少原子在链上；缩合聚合每接一节都要“付一笔水费”。

\begin{qiaoliang}
一根聚合物链到底有多长？来估一估。取一根摩尔质量约 \ud{140000}{g\;mol^{-1}} 的聚乙烯链：每个乙烯单元的质量是 \ud{28}{g\;mol^{-1}}，所以链上约有 $140000\div 28\approx 5000$ 个单元，也就是约 10000 个碳—碳单键。每个单键长约 \ud{0.154}{nm}，若把链拉成直线，长度约为 $10000\times 0.154\approx 1540$ 纳米，即大约 1.5 微米——考虑链其实是锯齿形，真实“伸直长度”约 1 微米出头，但量级不变。而这根链平时蜷成一团无规线团，直径只有几十纳米。换句话说：\textbf{一根链拉直后比它平时占的地方长几十倍}。你手指捏起的一小块塑料里，这样的链有 $10^{18}$ 根以上，彼此穿插缠绕——这就是那碗“面条”的真实密度。
\end{qiaoliang}

\section{调材料的四个旋钮}

化学家手里有四个“旋钮”，拧一拧，同一种原料能变出天差地别的材料。四个旋钮都建立在同一条规则上：链内是牢固的共价键，链间是可滑动的弱作用力。

\begin{itemize}[itemsep=2pt]
\item \textbf{链长}：链越长，缠得越死，材料越韧、越难熔融流动。二三十个碳的“迷你链”是蜡烛油，一捏就碎；几千个碳才是塑料袋；上万个碳的超高分子量聚乙烯能做防弹衣和人工关节。
\item \textbf{支化}：支链多的 LDPE 像一堆长满树杈的树枝，堆不紧密，链间作用力小，所以软、熔点低、透明；几乎无支链的 HDPE 像一捆筷子，排得整整齐齐，链间贴得紧，所以硬、结实、乳白色。
\item \textbf{结晶度}：链段如果能整齐排列，就会形成微小的晶区（第 24 章讲过晶体与玻璃的区别）。晶区像一颗颗“钉子”，把周围的链钉在原地——结晶度高的材料更硬、更强、更耐热，但光线穿过晶区边界时会散射，所以透明下降。PET 瓶子又透又韧，正是因为工厂把结晶度控制在恰到好处的低水平。
\item \textbf{交联}：在链与链之间用真正的化学键搭“桥”。桥很少时，链段还能滑、还能被拉直，但桥把链彼此拴住，一松手链就弹回原位——这是橡胶。桥很多时，整团链焊成一张三维大网，再也滑不动——这是电木插座、环氧树脂，加热到烧焦也不会熔。
\end{itemize}

最后一个旋钮直接划出了材料世界的两大阵营：\textbf{热塑性}材料（没有交联或极少，受热变软、冷却变硬，可以反复熔了再造）和\textbf{热固性}材料（交联成网，受热不熔，只能炭化分解）。饮料瓶、塑料袋是热塑性；轮胎、插座外壳是热固性。将来你会看到，这一字之差决定了哪种塑料能回收、哪种只能填埋。

橡胶是个值得单独说的案例。天然橡胶来自橡胶树的乳汁，单体是异戊二烯，长链天生蜷曲，所以生橡胶软塌塌、发黏、一拉就断。1839 年，美国人古德伊尔偶然把混了硫黄的生橡胶落在热炉子上，发现它变得干爽有弹性——硫原子在链与链之间搭起了少量的桥（“硫化”）。桥不多，链段仍能拉直；桥又足够，链不会彼此溜走。今天的轮胎用的就是这个思路，只不过配方早已精细得多。

还有一种“不像塑料的塑料”：硅胶。奶嘴、锅铲、烤箱垫、手机壳里的它，主链不是碳—碳，而是硅和氧交替相连（\chem{Si{-}O{-}Si{-}O}），硅上再挂两个小有机基团。硅—氧键更耐高温、更柔顺，所以硅胶能进烤箱也能进冰箱。这提醒我们：聚合物的主链不一定非得是碳。

再把乙烯玩出最后一个花样：把链上所有的氢原子都换成氟，得到聚四氟乙烯（PTFE，商品名“特氟龙”），就是不粘锅上那层黑亮的涂层。碳—氟键极牢，氟原子像一层严丝合缝的盔甲裹住碳链，别的分子既抓不住它、也攻不进去——食物粘不上，强酸强碱也奈何不了它。但它并非金刚不坏：空锅干烧到两三百度以上，链也会开始分解，所以不粘锅的正确用法从来不是“随便造”，而是别空烧、别用铁铲刮。

为什么橡皮筋室温下软、PET 瓶一灌开水就变形收缩？这里藏着一个温度开关：每条链都有一个“冻结温度”（玻璃化温度，第 24 章提过玻璃），低于它，链段冻得动不了，材料又硬又脆；高于它，链段开始扭动，材料就软。PET 的冻结温度大约在 $70\,^{\circ}\mathrm{C}$ 上下，开水一烫就越过了开关；橡胶的冻结温度低到零下六七十度，所以整个冬天它都是软的。

\section{现在，让符号登场}

讲了这么多链，终于可以写式子了。

加成聚合这样记：$n$ 个乙烯分子打开双键，接成一条链——

\[n\,\chem{CH_2{=}CH_2} \rightarrow \chem{{[CH_2{-}CH_2]}_{n}}\]

方括号里是\textbf{重复单元}，下标 $n$ 表示“这个单元在链上重复了 $n$ 次”。把乙烯的一个氢换成氯，得到氯乙烯，聚合后就是水管、地板革用的聚氯乙烯（PVC）：

\[\chem{{[CH_2{-}CHCl]}_{n}}\]

硬 PVC 做水管，掺入大量增塑剂（一种能钻进链间、让链更容易滑动的小分子）就变软 PVC，做电线外皮和桌垫。

缩合聚合只记一句话账：对苯二甲酸 $+$ 乙二醇 $\rightarrow$ PET $+$ 水，己二酸 $+$ 己二胺 $\rightarrow$ 尼龙 $+$ 水。每生成一个酯键或酰胺键，就有一分子 \chem{H_2O} 离队。

天然橡胶的重复单元来自异戊二烯，写作 \chem{CH_2{=}C(CH_3){-}CH{=}CH_2}，聚合后每个单元上仍留一个双键——正是这些残留的双键容易被氧气和臭氧攻击，橡皮筋晒久了才会发黏断裂。轮胎要掺碳黑和防老剂，就是在保护这些薄弱环节。

有一点必须说破：同一个瓶子里，各条链的 $n$ 并不相同，有的三千，有的五千。所以聚合物没有单一的“分子量”，只有一个平均值和一个分布——高分子材料天生就是混合物，这与你在第一册认识的纯净物（比如水、氯化钠）很不一样。

\begin{zhengju}
今天“分子可以有几千个原子”听起来天经地义，但在 1920 年代，这是化学界吵得最凶的问题之一。当时的主流看法是：橡胶、淀粉、纤维素这些材料是“胶体”——一群小分子靠某种神秘的“余价力”松散地抱成团，并不存在真正的大分子。理由听起来也挺硬：用 X 射线测这些材料，晶胞都很小，所以分子也应该小。

1920 年，德国化学家施陶丁格提出相反的假说：橡胶就是异戊二烯用普通共价键连成的长链，没有什么神秘力量。假说不能只靠嘴说，他设计了一个判决性实验：如果橡胶是靠双键间的“余价力”聚拢的小分子，那么用氢气把双键全部加成掉（氢化），团块就该散架变小；如果橡胶本来就是共价长链，双键没了，链依然还是长链。实验结果：氢化后的橡胶依然保持着高分子的行为。神秘力量说不攻自破——毕竟被拆掉的正是它假想的“胶水”。

接着是定量证据。施陶丁格发现，同种聚合物的溶液越稠（黏度越大），链就越长，黏度可以当“分子链的秤”用；用渗透压和后来的超速离心机测分子量，得到的结果是几千到几十万——比当时任何人愿意相信的分子都大两个数量级。争论一直持续到 1930 年代，直到美国杜邦公司的卡罗瑟斯用缩合聚合一步步合成出结构明确的尼龙（1935 年），把“人可以一节一节亲手搭出大分子”变成事实，胶体学派才真正退场。1940 年尼龙丝袜上市，第一天就卖出几十万双；1953 年，施陶丁格拿到诺贝尔化学奖。

这个故事的要点不在谁聪明，而在于：两种解释都能说通“材料又韧又黏”这一现象，分出胜负靠的是\textbf{设计一个让两种解释给出不同预言的实验}——氢化实验就是这样的分岔路口。
\end{zhengju}

\begin{shiyan}
\textbf{拉一拉，看分子链排队。}

材料：一个干净的聚乙烯塑料袋（超市购物袋，确认是 PE 材质）、一根橡皮筋。

步骤：从塑料袋上剪下一条宽约 2 厘米、长约 15 厘米的条带，两手捏住两端，缓慢匀速地拉，眼睛盯住条带中段。拉到一定程度后停下，观察条带外观的变化。再换橡皮筋，慢慢拉长，把拉长的橡皮筋轻轻贴在上嘴唇（那里对温度敏感），停留两三秒，然后松手。

观察点：塑料袋条被拉到某一处时，会突然出现一段\textbf{变细、变白}的“细脖子”，继续拉，细脖子会向两端蔓延——这叫“颈缩”，是混乱的链被强行捋直、重新排列的过程，一旦排过队的地方就再也缩不回去。橡皮筋则完全不同：拉长的瞬间你会感到它微微发热，松手后它能回到原长。

安全提示：本实验只用双手，不加热、不点火。\textbf{绝对不要}烧塑料观察现象——塑料燃烧会放出刺激性甚至有毒的气体，想看作废塑料制品的处理，请看回收站的分类，或者看老师的演示视频。剪下来的塑料条用完丢进可回收物桶。
\end{shiyan}

\section{这个模型在哪里会失手}

“链内强键、链间弱力、四个旋钮”是理解高分子的主干，但你手上的真实塑料制品，远比模型复杂。

第一，几乎没有商品塑料是“纯聚合物”。PVC 里有增塑剂，轮胎里有碳黑，矿泉水瓶里有防止加工时降解的稳定剂，彩色塑料里有颜料。配方属于工程，公开的资料往往只写主要成分——所以仅凭“这是 PET”并不能推断它的一切。

第二，分子量是个分布，结晶度也只是个平均值：同一块材料里，晶区和非晶区犬牙交错。模型能解释趋势（“支化多则软”），却算不出某一只具体瓶子的精确强度。

第三，“降解”这个词比听起来狡猾。塑料袋在太阳下变脆、碎裂，常常只是长链被打断成短链、整片碎成微粒，碳原子并没有回到自然界的循环里——从“看得见”变成“看不见”，不等于消失。标着“可降解”的塑料，多数需要在工业堆肥的高温高湿条件下才会真正分解，随手丢进草丛或海里，它和普通塑料一样能躺很多年。判断一种材料的环保账，要看从原料、制造、使用到废弃的整条链——这叫全生命周期评价，我们马上回来算这笔账。

\begin{wuqu}
“聚合物都是人工合成的塑料，天然的东西里没有聚合物。”——错得离谱，而且一错就错过半本书。淀粉和纤维素都是葡萄糖串成的聚合物（第 46 章细讲为什么同一种单体串出的两种链，一种你能消化、一种不能）；蛋白质是二十种氨基酸串成的聚合物（第 48 章）；连你身体里的遗传物质 DNA，也是四种核苷酸串成的超长高分子。人类发明塑料，满打满算一百多年；大自然玩“单体串长链”这套把戏，已经玩了三十多亿年。化学家真正发明的，只是用石油里的单体来串而已。
\end{wuqu}

\section{回到那只瓶子}

现在兑现开场的承诺。

矿泉水瓶：它是对苯二甲酸和乙二醇“牵手付水费”串出的 PET 链，每条约一两百个单元。轻，是因为链本身密度低、堆得又不算紧；摔不碎，是因为上万条缠结的链能通过滑动来吸收冲击，而不是像玻璃那样一裂到底；透明，是因为工厂把结晶度压得很低；而它又足够结实，是因为吹瓶时链被双向拉伸、排了队（和你的塑料袋实验里的“细脖子”同理）。它不烂，不是细菌不想吃，而是碳骨架上没有酶认得的接口，酯键的水解在室温下慢得几乎没有动静。

橡皮筋：它是聚异戊二烯长链加少量硫桥。能拉长，是蜷曲的链被拉直；能弹回，是硫桥拴着链不让它们溜走；晒久了发黏断裂，是残留的双键被氧气、臭氧和紫外线剪断，交联网络破了洞。

涤纶衣服：它和矿泉水瓶是同一种 PET，只不过被熔融后从细孔里挤出、拉成丝——所以标签上“涤纶”“聚酯纤维”“PET”指的是同一条链。快干，是因为这条链上能抓水的地方很少；遇火星熔出小洞，是热塑性链一遇高温就软化流动。

三样东西，一条逻辑：链内共价键决定“牢不牢”，链间作用力和交联决定“动不动”，而“动不动”几乎决定了你摸到的全部手感。

\section{接着往下挖：一只瓶子的身后事}

捏一捏饮料瓶底部，你会看到一个三角形里套着数字。那是塑料的“回收身份证”：

\begin{table}[htbp]
\centering
\begin{tabular}{clll}
\toprule
编号 & 材料 & 常见身影 & 回收去向 \\
\midrule
1 & PET（聚酯） & 饮料瓶、食用油瓶 & 回收成熟，可再纺成纤维 \\
2 & HDPE（高密度聚乙烯） & 牛奶桶、洗发水瓶 & 回收较成熟 \\
5 & PP（聚丙烯） & 外卖餐盒、瓶盖 & 部分地区回收 \\
6 & PS（聚苯乙烯） & 泡沫饭盒、一次性杯盖 & 回收困难 \\
7 & 其他（含 PC、亚克力等） & 水壶、镜片、混合材料 & 基本不回收 \\
\bottomrule
\end{tabular}
\end{table}

这张表背后是一个工程难题：回收要求把不同种塑料分开，因为混在一起的链熔点不同、互不相容，再生出来的材料性能一塌糊涂。热塑性塑料至少还能熔了再造，热固性的轮胎和电路板连这一关都过不了。于是大量塑料的去处是填埋、焚烧，或者——泄漏进环境，在日晒海浪中碎成直径小于 5 毫米的微塑料，出现在深海、土壤、甚至我们的血液里。微塑料对健康的影响目前仍在研究之中：实验室里高剂量的确能伤害细胞和动物，但现实中人们吃进去的剂量有多大、会造成什么后果，证据还不够强——这是一个“值得警惕、但还不能下结论”的问题，记得用“剂量决定毒性”的眼光看待新闻标题。

那干脆不用塑料、全换回纸和棉布行不行？这就要算全生命周期的账了：纸袋更重、更耗水和树，一只棉布袋的碳排放账单远高于一只塑料袋，一些研究估算它要重复使用几十到上百次，气候账才能追平。答案不是“禁掉哪种材料”，而是：让每条链被用得足够久、回收得足够干净。

最后留一个悬念。这一章讲的是人类用两种串法造长链；可是大自然早就在做更惊人的事：它只用葡萄糖一种单体，串出了你能吃的淀粉和你不能吃的纤维素；它用二十种氨基酸按精确顺序串链，让链自己折叠成会干活的机器。链和链的差距，怎么就会大到“燃料”和“机器”的区别？从第 45 章开始，我们要走进大自然的聚合物工厂——在那里，你会重新见到这一章的每一个概念：单体、缩合、链间作用力、交联，只是工厂的账本，换成了生命自己的算法。

\begin{tiaozhan}
拉橡皮筋时它为什么会发热？答案违反直觉：橡胶的弹性主要不是“键被拉长后的回弹”，而是\textbf{熵}在驱动（第 28 章）。蜷曲的链有巨量的微观排布方式，熵高；被拉直的链排布方式骤减，熵低。松手后，链在热运动的撞击下随机乱窜，走向高熵的蜷曲态——宏观上就是收缩回弹。拉伸时链被迫有序，热运动能量以热的形式释放出来，所以你会感到温热；反过来，给一根绷紧的橡皮筋加热，它会收缩变短，而不是像大多数固体那样膨胀。这还预言了一件可以检验的事：把橡皮筋挂在零下几十度的低温里，它会越过冻结温度，变成一掰就断的“玻璃”。热力学的统计规律，就这样藏在你的笔袋里。
\end{tiaozhan}

\section*{练一练}

\begin{sikaoti}
\item 画两张物质流图：第一张画 $n$ 个乙烯聚成聚乙烯，标出双键在哪里打开、副产物有没有；第二张画对苯二甲酸和乙二醇缩聚成 PET，标出水分子在哪一步离开体系。比较两张图的“原子账”。
\item 只用“链长、支化、交联、结晶度”四个旋钮，解释：为什么牛奶桶（HDPE）乳白色、硬挺，而塑料袋（LDPE）透明、柔软？两者的单体明明一样。
\item 估算：某聚氯乙烯链平均含 1000 个重复单元，每个单元的摩尔质量约 \ud{62.5}{g\;mol^{-1}}。这条链的摩尔质量大约是多少？再按每个单元沿链方向约 0.25 纳米估算它拉直后的长度，并与它蜷成线团后的尺寸（几十纳米）比较。
\item 饮料瓶能熔化后重新纺丝，轮胎却只能粉碎后铺路。画出两者链结构的示意图（直线表示链、短横表示交联桥），用图说明为什么热固性材料“熔不了”。
\item 淀粉和纤维素都由葡萄糖串成，人类能消化淀粉却不能消化纤维素。结合本章“酶认接口”的思路，猜一猜：决定你能不能消化一条天然高分子链的关键因素可能是什么？（第 46 章揭晓。）
\end{sikaoti}
